\documentclass[11pt]{article}
\usepackage[margin=1.1in]{geometry}
\usepackage{amsmath,amssymb,amsthm}
\usepackage{booktabs}
\usepackage{array}
\usepackage{longtable}
\usepackage{microtype}
\usepackage{xcolor}
\usepackage[colorlinks=true,linkcolor=blue!50!black,urlcolor=blue!50!black]{hyperref}

\newtheorem{lemma}{Lemma}
\newtheorem{corollary}{Corollary}
\newcommand{\aggr}{\rho_{\mathrm{agg}}}
\newcommand{\rfill}{\rho_{\mathrm{fill}}}
\newcommand{\bMat}{B}
\newcommand{\code}[1]{\texttt{#1}}

\title{\texttt{sdpa-dd-omp}: threading a double-double SDP solver,\\
and making its threaded results reproducible}
\author{\texttt{sdpa-dd-omp} --- technical record}
\date{2026-08-25}

\begin{document}
\maketitle

\begin{abstract}
\noindent
\code{sdpa-dd} solves semidefinite programs in double-double arithmetic --- roughly $32$
significant decimal digits, carried as an unevaluated sum of two IEEE doubles. It builds with
OpenMP, but threading it as shipped is not merely unprofitable, it is unsound: on small problems
more cores make it monotonically slower ($24.5\times$ slower on $24$ cores than on one, on a
user's $m{=}16$ input), and on any problem its threaded runs are not reproducible.

This document records what was changed and why. Five things are worth a reader's attention.
First, the dominant cost is the Schur-complement build and factorisation, and on large sparse
problems upstream threads \emph{neither} the sparse factorisation nor its assembly, so
twenty-four cores buy it $4.9\%$. Second, threading a multiprecision solver is mostly a
\emph{reproducibility} problem: this fork's assembled Schur complement and its Cholesky factor are
bit-identical at any thread count \emph{by construction}, and that is asserted against the
structures themselves rather than against printed digits. Third, one of this fork's own released
versions shipped a data race, and the account of how it was missed is more useful than the fix.
Fourth, the rule choosing between a dense and a sparse factorisation of the Schur complement was
mis-tuned in a way a retune would not have found: one of its four gates is not a heuristic at all
but a provable consequence of another. Fifth, on \code{aarch64} the compiler's default
floating-point contraction silently changes double-double arithmetic, so results differed between
platforms until contraction was pinned off.
\end{abstract}

\tableofcontents

\newpage

\section{The solver, and where the time goes}

At each interior-point iteration the solver forms the $m \times m$ Schur complement $\bMat$ (the
``\code{bMat}''), factorises it, and performs a forward and a backward triangular solve against
the factor. In double-double arithmetic every scalar operation is an inlined sequence of some ten
to twenty IEEE operations with strict ordering requirements, so the constant factors a
double-precision solver hides become the whole story.

Two regimes matter, and they have different bottlenecks.

\paragraph{Small and medium problems ($m \lesssim 10^3$).} $\bMat$ is dense, and the cost is the
constraint loop that builds it plus a dense Cholesky. This is the regime SDPLIB covers, and it is
the regime upstream's own threading was written for.

\paragraph{Large sparse problems ($m \sim 10^3$--$10^4$).} $\bMat$ has a sparsity structure worth
exploiting, and the solver can store and factor it sparsely. Here the cost splits between the
\emph{sparse assembly} --- scattering constraint-pair contributions into the stored nonzeros ---
and the \emph{sparse Cholesky}. \textbf{Upstream threads neither.} This is why twenty-four cores
buy upstream $4.9\%$ on \code{dE4} ($m{=}7401$): the problem routes sparse, and both of the
regions that dominate it are serial code. It is also why the SDPLIB tables understate this fork:
every SDPLIB problem takes a dense $\bMat$ and never enters the path where the difference lives.

A methodological note, recorded because the discarded hypothesis is part of the result: the first
optimisation proposed for the sibling \code{sdpa-gmp} port was a per-block parallel Cholesky over
the SDP blocks. Profiling assigned $66$--$80\%$ of runtime to the $\bMat$ build, and the proposed
region did not appear in the profile at all. It was abandoned on the evidence.

\section{Kernel and threshold defects in the inherited code}

\subsection{The lost zero-skip in \code{Rgemm}}

Netlib \code{dgemm} guards its inner update with a test for a zero multiplier. In
\code{mplapack}~2.0.1 that guard was dropped from the OpenMP variants. In double precision the
guard is a minor optimisation; in double-double it is not, because ``multiply by zero'' still pays
a full multiprecision call. Restoring it in \code{mplapack/Rgemm\_NN\_omp.cpp} and
\code{Rgemm\_NT\_omp.cpp} recovers up to $2\times$ on the serial path.

\subsection{The bundled QD builds without hardware FMA on aarch64}
\label{sec:qdfma}

QD's own \code{configure} runs its FMA probe only under \code{case \$host in powerpc*-*-*)}. The
default \code{--enable-fma=auto} therefore resolves silently to \emph{none} everywhere else:
\code{QD\_FMA}/\code{QD\_FMS} stay undefined, and \code{two\_prod} --- the most-executed primitive
in the whole solver --- compiles the Dekker two-way split instead of
$p = a b;\ \mathrm{err} = \mathrm{fma}(a, b, -p)$.

On Apple Silicon that is \textbf{72 instructions against 8} (\code{dd\_real::operator*}: 71 against
13), the difference being \code{fmul}+\code{fnmsub} against a sequence containing two conditional
branches for \code{split()}'s $2^{996}$ overflow guard. Enabling it moves the solver's main loop by
\textbf{$1.59\times$ in aggregate} at one thread --- per problem $1.21$--$1.63\times$ (\code{truss6}
$1.21$, \code{arch0} $1.38$, \code{mcp250-1} $1.53$, \code{gpp250-1} $1.61$, \code{control7} $1.62$,
\code{theta3} $1.63$) --- against a repeat spread of $0.1$--$2.8\%$, so the smallest ratio exceeds
the largest spread by about $4\times$ and no row is marginal.

\paragraph{Bit-identical, and measured rather than assumed.} Sixteen SDPLIB problems were hashed
over the complete solution section: reference-versus-candidate at 8 threads $0/16$, candidate 1
thread versus 8 threads $0/16$, reference-versus-candidate at 1 thread $0/16$, and both arms under
\code{-ffp-contract=on} $0/16$. Iteration counts were identical across all $96$ runs. The comparison
is proven \emph{sensitive} rather than merely quiet: perturbing \code{lambdaStar} by one part in
$10^{13}$ with the same binary moved $17$ fields across $6$ problems. A scalar cross-check over
$\sim$$106$M random pairs found the product and the error term agreeing exactly across the solver
band, the normal band, \code{two\_sqr}, near-overflow and huge-times-tiny; they diverge only for
deep subnormals, where the exact error term is below \code{DBL\_MIN} and unrepresentable.

\paragraph{Strictly more correct, too.} Near \code{DBL\_MAX} the Dekker split overflows and returns
\code{inf}/\code{nan} for a finite product --- $a = 10^{-300}$, $b = \code{DBL\_MAX}$ gives
$\mathrm{err} = \infty$ under the split and $9.84\times10^{-9}$ under FMA. That is a latent defect
in the non-FMA builds, removed here on aarch64.

\paragraph{The aarch64 gate is deliberate, not an oversight.} On aarch64 \code{fmadd} is baseline
ISA, so \code{\_\_builtin\_fma} is a single instruction. On x86-64 \emph{without} \code{-mfma}, GCC
lowers \code{\_\_builtin\_fma} to a libm \emph{call}, which is far slower than the split --- and
this project builds generic x86-64 on purpose (no \code{-march}, no \code{-mfma}) so binaries stay
bit-identical across AMD and Intel nodes. On every non-aarch64 host the argument handed to QD's
\code{configure} is unchanged, and x86-64 was verified untouched at seven levels, including $74/74$
objects identical after stripping, \code{.text} byte-identical, and all $60$ published state-hash
cells regenerated and matched. \code{--enable-qd-fma=auto|gnu|no} overrides the host gate, and
exists so that ``the same build with QD's explicit FMA turned off'' is configurable without
hand-editing --- hand-editing does \emph{not} invalidate \code{external/qd}'s configure stamp, and
would silently reuse the previous library.

One inherited trap was fixed alongside it: \code{work/.qd\_configure\_done} depended only on the
patch stamp, so an existing build tree kept a stamp made with the old arguments and silently
produced a non-FMA library. Anyone benchmarking in an existing tree would have measured the
reference while believing they measured the candidate.

\paragraph{Known gap, deliberately not closed here.} \code{install.sh} short-circuits to a system QD
on Linux when \code{/usr/include/qd/qd\_real.h} exists, so an aarch64 Linux box with a distribution
\code{libqd-dev} never reaches this flag. Fixing that would also change which QD x86-64 Linux picks
up, and so needs its own decision.

\subsection{OpenMP regions entered regardless of problem size}

Upstream enters its parallel regions without asking whether there is enough work to repay the
fork/join. It contains a threshold mechanism, but the mechanism is disabled behind \code{if (0)},
and it cannot simply be switched on: the \code{\_ref} bodies it dispatches to are absent and the
tree would not link.

The consequence is not a small inefficiency. On a user's \code{max\_custom\_scaled.dat-s}
($m{=}16$), upstream degrades \emph{monotonically} at every thread count, ending $24.5\times$
slower on $24$ cores than on one. On SDPLIB \code{control1} it is $11\times$ slower at $24$
threads than at $1$; on \code{truss5}, $3.3\times$ slower. This fork gates entry on work, matrix
width and nesting depth, and runs \code{Rdot} serially; \code{control1} is then flat to four
decimal places across the whole thread ladder. The gain on small problems reaches $13\times$, and
it also removes a source of nondeterminism, since a reduction whose combining order followed
thread scheduling was the mechanism behind upstream's varying iteration counts.

\subsection{Timers that reported worker-seconds as elapsed time}

The reduced per-formula timers summed time across threads and printed the sum where elapsed time
was expected, so a threaded run appeared to spend more time in a region the more it was
parallelised. They are now reported as worker-seconds, labelled as such.

\section{Reading high-precision input}

This defect is not a performance matter and is the one most likely to have silently corrupted
someone's results.

QD's decimal-to-\code{dd\_real} converter overflows internally once the mantissa reaches $309$ or
more digits, and \textbf{returns success while producing NaN}. A $600$-digit spelling of an
ordinary-magnitude number therefore entered the SDP matrices as NaN and surfaced later as a bogus
\code{cholesky miss condition} at iteration~$0$ --- the solver reporting an infeasible problem
when the problem was fine.

The isolation is clean. On a user's $600$-digit input, upstream terminates after
\textbf{0 iterations} with \code{objValPrimal = -0.0}; this fork reaches \code{pdOPT} in
\textbf{43 iterations}. Truncating that same input to \textbf{45 significant digits} --- beyond
double-double's own $\sim$$32$-digit resolution, and far below the converter's $\sim$$309$-digit
limit --- makes \emph{upstream} solve it correctly. Token length is the only thing that changed.
Nor is the truncation doing arithmetic work: this fork's answers on the $600$-, $45$- and
$30$-digit inputs agree to ten significant digits, and a difference that small cannot turn a NaN
into a correct solve.

The reader now validates full decimal syntax, keeps the original conversion path for mantissas up
to $308$ digits so ordinary inputs stay bit-identical, normalises longer mantissas with the
exponent preserved, and checks parser status \emph{and} finiteness rather than letting a NaN reach
the solver.

\section{Floating-point contraction, and why this fork pins it off}

\subsection{The mechanism}

GCC and clang default to \code{-ffp-contract=fast} for C++. Double-double arithmetic is
header-inline, so \emph{the compiler} decides whether the compensation terms of the
double-double algorithms get folded into a fused multiply-add. Folding there is not an identity.
The whole point of a \code{two\_sum} is to compute the rounding error that the primary operation
discarded; an FMA that keeps the intermediate product exact deletes the very quantity being
measured, and the double-double invariant is lost.

The platform dependence follows immediately. On baseline \code{x86-64} there is no FMA for the
compiler to fold into, so those binaries behave as if contraction were off. On \code{aarch64},
\code{fmadd} is baseline ISA, so they do not. The two platforms then run different arithmetic on
the same source.

\subsection{The fix, and its cost}

This fork builds with \code{-ffp-contract=off} \textbf{by default}, applied to every SDPA and
\code{mplapack} translation unit \emph{and} to the bundled QD. With it, an M1 build matches the
published \code{x86-64} reference on \textbf{20 of 20} SDPLIB problems --- iteration counts and
full-precision objectives alike --- where the unpinned build matches on none at the state-hash
level.

Hardware FMA is not lost. \code{-ffp-contract=off} does not disable an explicit
\code{\_\_builtin\_fma}, so QD's \code{two\_prod} keeps its single \code{fmsub} and the
\code{aarch64} speedup of \S\ref{sec:qdfma} is retained in full --- verified at the instruction
level, not assumed. The cost is $\sim$$1\%$ per iteration at one thread
and $\sim$$1.7\%$ at eight. On \code{x86-64} the output is unchanged and so is the object code:
all $76$ translation units compile to byte-identical \code{.text} with and without the flag.

\code{--enable-fp-contract=fast} recovers the old, platform-dependent behaviour. Iteration counts
published for \code{arm64} \emph{before} this became the default are counts of the unpinned
solver; the per-problem deltas are in \code{bench/b4\_mac\_rebaseline.tsv}.

\section{What was threaded, and what it costs correctness}

Four regions are threaded here that are serial or ineffectively threaded upstream.

\subsection{The \texorpdfstring{$\bMat$}{bMat} constraint loop}

The dense-route Schur-complement build. This was the largest single gain in the original port and
is where the SDPLIB improvements come from.

\subsection{The sparse Schur-complement Cholesky}

\textbf{Bit-identical at any thread count by construction}, and this is a property of the loop's
dependence structure rather than of the arithmetic --- so it does not depend on the working
precision, nor on the values of the work gates. For a fixed pivot, each \code{k1} owns exactly one
destination row; every read comes from the pivot row; and a barrier per pivot preserves pivot
order. No sum is reassociated, because no destination is written by two workers.

\subsection{The sparse Schur-complement assembly}

Parallel over $i$-groups, serial over blocks. Each worker holds private \code{work1}/\code{work2}
scratch matrices, so concurrent groups cannot race through shared scratch, and the accumulation
into $\bMat$ is disjoint \emph{within} a block. That last qualification is load-bearing and is the
subject of \S\ref{sec:race}.

\subsection{Cholesky panel kernels and the \texorpdfstring{$X$/$Z$}{X/Z} inverse-Cholesky triangulars}

Threaded behind measured work gates.

\subsection{The triangular solves: forward yes, backward no}

The forward substitution is a \emph{scatter}: within a row the writes go to distinct destinations,
so there is no reassociation and the threaded result is bit-identical by construction. The sibling
\code{sdpa-gmp-omp} threads it on that basis.

The backward substitution is a \emph{dot product}, and splitting a dot product across workers
reassociates the sum. In \code{sdpa-gmp} this is recoverable: GMP carries a runtime precision, so
the accumulator can be widened with \code{mpf\_set\_prec(acc, 2*prec)} and the partial sums
combined without loss. \code{dd\_real} has no runtime precision to double --- the type \emph{is}
its precision --- so the same remedy does not exist here. In \code{sdpa-gmp-omp} the backward pass
is therefore opt-in behind \code{SDPA\_SOLVE\_BACKWARD} and off by default; in this fork it is not
threaded. Trading bit-reproducibility for a fraction of one phase is not a trade this fork makes
silently.

\section{A data race in the threaded assembly, and its fix}
\label{sec:race}

\subsection{What shipped}

The threaded sparse assembly's worksharing loop carried \code{nowait}, justified by the observation
that the $\bMat$ writes are \code{+=} into distinct slots \emph{within a block}. That observation
is true and insufficient. \code{nowait} lets one worker enter block $l+1$ while another is still in
block $l$, and \emph{across} blocks the slots are not distinct: a constraint pair that is nonzero
in two blocks accumulates into the same entry from both.

The consequence, on SDPLIB \code{truss6} at default settings, where the correct objective is
\code{-9.0100140477088098e+02}. Five consecutive runs of \code{v7.1.3-omp.2} returned

\begin{center}
\small
\begin{tabular}{@{}lll@{}}
\toprule
\code{2.1611901436772943e+06} & \code{9.7205414006274572e+04} & \code{-9.0093934416427931e+02} \\
\code{1.6632246196208220e+04} & \code{9.8204796932176279e+05} & \\
\bottomrule
\end{tabular}
\end{center}

The third value is the instructive one: it is wrong in the \emph{fifth significant digit}, not
absurd. A tolerance check against a reference would have accepted it. On one host, three of five
wrong values agreed with each other to four digits, so a small-sample spread test would have
called the run stable.

\subsection{Why nothing caught it}

Three reasons, none of them ``we did not test''.

\begin{enumerate}
\item \textbf{The CI fixture was block-diagonal by construction.} No Schur entry is ever written
      twice in it, so the race had nothing to land on. Every stream-identity assertion ran on that
      fixture, and every one of them passed.
\item \textbf{The large problems were not immune, and were believed to be.} Both \code{dE3} and
      \code{dE4} report \code{blocks\_disjoint=0}. The race fired on \code{dE4} \emph{once in
      seventeen} threaded runs. They escaped notice because a rare race compared \emph{once} is a
      race not tested for.
\item \textbf{The guarantee was stated in a form that misdirected every check.} It was always
      written as ``bit-identical at any thread count'', which points every test at comparisons
      \emph{between} thread counts. A race that moves both sides passes those comparisons. The
      check that catches it is repetition at \emph{one} thread count, and nobody had written it.
\end{enumerate}

\subsection{The fix}

The barrier is taken, but only where the finished aggregate map \emph{proves} it is needed. One
setup pass asks whether any Schur entry is written by more than one SDP block:

\begin{quote}
\code{if (blocks\_disjoint)} $\rightarrow$ \code{\#pragma omp for schedule(dynamic) nowait} \\
\code{else} $\rightarrow$ \code{\#pragma omp for schedule(dynamic)}
\end{quote}

with a single extracted \code{Newton::assemble\_group()} body shared by both worksharing forms, so
the two paths cannot drift apart. Block-diagonal problems keep \code{nowait}; cross-block problems
take the barrier. \emph{The fast path is taken because the map was checked} --- which is precisely
the discipline whose absence caused the defect.

The disjointness pass itself uses two bitsets over the stored Schur entries, \code{seen\_any} and
\code{seen\_blk}, with \code{seen\_blk} cleared by re-walking the block's own locations rather than
by maintaining a scratch list. That is $O(\text{pairs})$ and two bits per stored entry, replacing
an earlier one-\code{int}-per-entry map that was $16\times$ larger.

Cost, on the four-iteration budget protocol: \code{dE4} $+0.7\%$, \code{dE3} via the fill route
$+2.2\%$, and nothing measurable in memory --- \code{dE4} is $382.0$~MB against $382.2$~MB before
the fix.

\subsection{The regression fixture, and why its first version was worthless}

\code{tests/gen\_crossblock\_fixture.py} generates an $m{=}480$ problem that routes sparse, threads,
and shares Schur entries between blocks. Its first version used a block stride of $7$, and the
\emph{racing} binary passed it $6/6$: blocks seven apart are never in flight together. Adjacency is
load-bearing, and the stride must be~$1$.

CI now asserts three things: that repeated runs at \emph{one} thread count agree; that the threaded
assembly stream equals the serial one; and that the overlap policy is the computed one \emph{in
both directions}, so a fixture that stopped exercising the cross-block case would fail rather than
pass vacuously.

\section{The factorisation choice}

\subsection{Four gates}

The chooser decides whether $\bMat$ is stored and factored densely or sparsely. It is a cascade of
four gates; a problem goes dense as soon as any gate rejects.

\begin{enumerate}
\item a coarse structural screen, which decides the great majority of problems;
\item a screen on the number of blocks contributing, against $c \cdot m$;
\item a pre-screen on the \textbf{aggregate} sparsity pattern of $\bMat$, against $F \cdot m^2$;
\item the real test, on the \textbf{ordered fill} of the symbolic Cholesky factor, against
      $F \cdot m^2$.
\end{enumerate}

Write $\aggr$ for the density of the aggregate pattern and $\rfill$ for the density of the ordered
fill, both as fractions of $m^2$.

\subsection{One constant doing two unrelated jobs}

Gates 2 and 3 were both driven by a single constant, so neither could be retuned without moving the
other. Splitting it frees gate~3 --- and once gate~3 is free, it stops being a tuned threshold at
all.

\subsection{Gate 3 is not a heuristic}

\begin{lemma}
Symbolic Cholesky factorisation only \emph{adds} entries: every nonzero of the aggregate pattern is
a nonzero of the factor's pattern. Hence $\rfill \ge \aggr$ always.
\end{lemma}

\begin{corollary}
$\aggr > F \implies \rfill > F$. So whenever gate~3 rejects at cutoff $F$, gate~4 would also have
rejected at cutoff $F$.
\end{corollary}

Gate~3 at $F$ is therefore a \emph{provable early exit} from gate~4, not an independent heuristic:
it can only skip work, never change the verdict. And that makes $F$ the \textbf{largest sound
cutoff} for gate~3. Anything smaller rejects problems gate~4 would have accepted.

\subsection{What the old value was, and why it was wrong}

Gate~3 rejected at $0.25 \cdot m^2$ while gate~4 rejected at $0.40 \cdot m^2$. Gate~3 was thus the
\emph{stricter} test, and every problem with $\aggr \in (0.25, 0.40]$ was sent dense
\textbf{without gate~4 ever running}. The cheap pre-screen overruled the real test.

The correction is to demote gate~3 to gate~4's own constant. \textbf{One tunable, $F = 0.40$,
unchanged, now replaces two.} It is worth being precise about what this is not: $0.25$ was not a
value tuned on evidence that later went stale. It was conservatism with nothing behind it, and the
soundness argument above shows there was never a regime in which it could pay.

\subsection{Who this affects}

Route decisions depend only on sparsity structure, so a census transfers across precisions and
across the sibling forks.

\paragraph{SDPLIB: nobody.} Of $92$ problems, \textbf{none} change route. $84$ are decided at
gate~1, $6$ at gate~2, and the single problem that reaches gate~3 --- \code{truss5} --- has
$\aggr = \rfill = 1.0$ and goes dense under both rules.

\paragraph{Bootstrap problems: most of them.} On a $221$-instance census, \textbf{167 instances
across seven structural classes} change dense~$\rightarrow$~sparse. Every one of those $167$ has
$\aggr \in [0.257, 0.362]$ and $\rfill \in [0.259, 0.367]$: above the old $0.25$, and
\textbf{not one above $0.40$}. Gate~4 said sparse on all of them.

\subsection{Measured effect}

One representative of \emph{each} of the seven structural classes was timed --- seven timed inputs,
one per class, not $167$. The classes are complete coverage of the switch set; the timings are
seven observations.

\begin{center}\small
\begin{tabular}{rrrrr}
\toprule
$m$ & \code{legacy} (dense) & \code{fill} (sparse) & faster & less memory \\
\midrule
$2{,}439$  & $1.92$~s / $116$~MB    & $0.61$~s / $50$~MB     & $3.17\times$ & $2.30\times$ \\
$4{,}489$  & $9.22$~s / $392$~MB    & $3.59$~s / $213$~MB    & $2.57\times$ & $1.84\times$ \\
$5{,}278$  & $14.50$~s / $519$~MB   & $3.88$~s / $234$~MB    & $3.74\times$ & $2.22\times$ \\
$6{,}067$  & $22.97$~s / $671$~MB   & $4.50$~s / $277$~MB    & $5.11\times$ & $2.42\times$ \\
$8{,}359$  & $66.20$~s / $1{,}289$~MB & $15.13$~s / $587$~MB & $4.38\times$ & $2.20\times$ \\
$10{,}614$ & $150.47$~s / $2{,}032$~MB & $22.36$~s / $831$~MB & $6.73\times$ & $2.45\times$ \\
$11{,}227$ & $175.81$~s / $2{,}358$~MB & $47.94$~s / $1{,}183$~MB & $3.67\times$ & $1.99\times$ \\
\bottomrule
\end{tabular}
\end{center}

Raw rows: \code{bench/dd-port3-2026-08-24/dd\_fill\_seven\_structures.tsv}.

\subsection{What this evidence does not establish, and where the rest of it comes from}

The table above is a \textbf{fixed four-iteration, per-iteration} experiment. \emph{No
time-to-solution comparison was performed on \code{dd}}, and none can be: these instances do not
converge at double-double precision under either tolerance tested (the shipped
\code{epsilonStar}${}={}$\code{1e-30} and a relaxed \code{1e-20}). A looser parameter choice is a
different experiment and nothing here rules it out.

The other half of the case comes from the sibling fork. \code{sdpa-gmp-omp} promoted the same
policy on 2026-08-18 after five review rounds, on the \emph{same seven structures}, measured
\textbf{at full convergence} ($m{=}6067$: dense $2708$~s against sparse $310$~s), on two
architectures, at $256$ and $512$ bits, at nine labelled thread points, with no reversal anywhere
--- and at \emph{one} thread, where sparse still won $2.60$--$3.06\times$, which shows the $0.25$
was already mis-routing in the serial regime it was calibrated for. Where the dense arm converged,
phases and iteration counts were equal per pair and the solution vectors agreed to
$6.3\mathrm{e}{-}38$.

\subsection{Not ``sparse always wins''}

\code{truss5} has an ordered fill of $1.0$ --- a completely dense factor --- and dense is genuinely
about $1.4\times$ faster there. Both the old and the new rule route it dense. That is the point of
testing the fill rather than assuming an answer.

\subsection{What the change costs, stated plainly}

\code{SDPA\_BMAT\_MODE} is \textbf{result-changing by design}. Dense and sparse are different
factorisation routes and follow different iterate trajectories, so results can differ in the last
digits. Across a $48$-cell harness, $30$ cells differ between routes, while every cell that
\emph{converges} agrees on the objective. Anyone needing bit-reproducibility against an earlier run
must keep the mode fixed; \code{SDPA\_BMAT\_MODE=legacy} reproduces the pre-2026-08-24 chooser
exactly.

\section{Reproducibility: established, not asserted}

\subsection{What is guaranteed}

\begin{itemize}
\item The assembled Schur complement is \textbf{bit-identical} at any thread count, and across
      repeated runs at one thread count.
\item The sparse Cholesky factor is \textbf{bit-identical} at any thread count, by the dependence
      argument of \S5.2.
\item A checksum over every scalar of the final \code{xMat}/\code{yVec}/\code{zMat} is identical at
      $1/2/4/8/24$ threads.
\item Results are identical between \code{arm64} and \code{x86-64}, given the contraction pin of
      \S4.
\end{itemize}

Unpatched upstream has none of these. Threaded, its iteration counts and even its objectives vary
run to run: \code{control1} ran $69/74/88$ iterations across three repeats at $24$ threads on one
host and $69/82/109$ at $8$ threads on another, where this fork ran $71/71/71$ on both.

\subsection{Why the printed solution is the wrong thing to compare}

A \code{dd\_real} carries about $32$ significant digits; the printed result fields carry $17$. A
factor can therefore change without any printed field moving. Every reproducibility claim here is
made against the \emph{structures} --- the assembled matrix, the finished factor --- and not
against printed digits.

\subsection{The oracles}

\begin{center}\small
\begin{tabular}{@{}p{5.6cm}p{7.6cm}@{}}
\toprule
variable & what it produces \\
\midrule
\code{SDPA\_SPCHOL\_DIGEST}          & fingerprint of the finished factor: records, bytes, 64-bit FNV-1a \\
\code{SDPA\_SPCHOL\_DIGEST\_DUMP=\emph{file}} & the factor's canonical stream, appended \\
\code{SDPA\_BMAT\_ASM\_DIGEST}       & the same fingerprint, for the assembled Schur complement \\
\code{SDPA\_BMAT\_ASM\_DUMP=\emph{file}} & the assembly's canonical stream \\
\bottomrule
\end{tabular}
\end{center}

The dumps are the proof-grade comparison: run twice, \code{cmp} the files. A fingerprint is strong
evidence, not proof, and is described that way wherever it is printed. Both streams carry a length
frame before every variable-length field, so two different structures cannot serialise to the same
bytes, and the two streams use different tags so an assembly stream can never compare equal to a
factor stream.

One operational caveat that cost a full sweep: the dump files are opened in append mode
(\code{"ab"}), so a comparison harness must \code{rm -f} them before each run or every problem will
appear to race.

\subsection{Negative controls, and why they are not optional}

A comparison that cannot fail is worse than no comparison, because it reports success. Each oracle
therefore has a hook that injects the smallest change it must detect, and CI asserts that the hook
actually changes what it claims to change. These require \code{-DSDPA\_SPCHOL\_TEST\_HOOKS} at
build time; a release build \emph{refuses} them rather than ignoring them, because a knob that
silently does nothing misleads whoever set it.

\begin{center}\small
\begin{tabular}{@{}p{6.1cm}p{7.1cm}@{}}
\toprule
hook & what it injects \\
\midrule
\code{SDPA\_SPCHOL\_MUTATE=1}    & one ulp into the finished factor \\
\code{SDPA\_BMAT\_ASM\_MUTATE=1} & one ulp into the assembled matrix \\
\code{SDPA\_BMAT\_TEST\_BREAK\_INVARIANT=1} & a violation of the $\rfill \ge \aggr$ invariant \\
\code{SDPA\_BMAT\_TEST\_F2\_STALE=1} & the worst-case stale-$G$ read a since-fixed defect could produce \\
\code{SDPA\_SPCHOL\_TEAM\_OVERRIDE=\emph{n}} & asks the runtime for $n$ threads, leaving the gates' team alone \\
\code{SDPA\_SPCHOL\_FAIL\_AT=\emph{i}} & declares pivot $i$ indefinite \\
\bottomrule
\end{tabular}
\end{center}

The last two are a different idea from the first four: not negative controls for an oracle, but the
only way to execute two branches that no input reaches on purpose.
\code{TEAM\_OVERRIDE=1} is the only route into the ``requested a team and received one thread''
fallback --- the dispatcher's own \code{team < 2} test exits earlier and by a different path, so
without the hook that branch is dead to every test --- and CI additionally requires the fallback's
factor to be byte-identical to the threaded one. \code{FAIL\_AT} makes the failure return from a
\emph{partly-threaded} factorisation executable, and CI requires that threaded pivots ran first, or
the test would prove nothing about that path. When it fires, the diagnostic says the failure was
\textbf{declared}, not measured.

\subsection{Route fixtures}

\code{tests/gen\_route\_fixture.py} generates eight cases straddling the gates and asserts the exact
route taken under \code{auto}, \code{fill} and with the variable unset. One case, \code{wide}, is
policy-divergent, and CI fails if it stops diverging --- because the assertion that unset behaves
like \code{auto} is vacuous unless at least one case can tell the policies apart.

\section{Environment variables}

Two rules hold for every variable below, and are tested in CI rather than merely stated.

\begin{enumerate}
\item \textbf{A typo is refused, never treated as the default.} \code{SDPA\_SPCHOL\_MODE=Auto} or
      \code{SDPA\_DD\_MIN\_SPCHOL\_WORK=-5} stops the run with a diagnostic. Silently falling back
      would let a mistyped variable select the very path the caller was trying to avoid --- and
      \code{-5} in particular used to be \emph{accepted}, wrapping to
      $18446744073709551611$, i.e.\ ``never parallel'' while looking configured.
\item \textbf{Validation happens on every route.} A knob is parsed once per process at a point both
      the dense and the sparse route reach, so a dense-route problem refuses a malformed value it
      would never otherwise consult. This fork has been bitten five times by checks placed where
      only one route could reach them.
\end{enumerate}

\subsection{Choosing how the Schur complement is stored and factored}

\begin{center}\small
\begin{tabular}{@{}p{5.1cm}p{5.6cm}l@{}}
\toprule
variable & values & changes results? \\
\midrule
\code{SDPA\_BMAT\_MODE} & \code{auto} (default) $\cdot$ \code{fill} (synonym) $\cdot$
  \code{legacy} $\cdot$ \code{dense} $\cdot$ \code{sparse} & \textbf{yes} \\
\code{SDPA\_BMAT\_MAX\_GB} & positive number of GB & no \\
\code{SDPA\_BMAT\_LOG} & \code{1} on, \code{0} off & no \\
\bottomrule
\end{tabular}
\end{center}

\code{auto} is the current default policy, which since \textbf{2026-08-24} is the fill-derived rule
of \S7; \code{fill} is retained as an explicit synonym so opt-in-era scripts keep working;
\code{legacy} restores the pre-2026-08-24 chooser exactly. \code{dense} and \code{sparse} force a
route. All values are strictly parsed. \code{SDPA\_BMAT\_LOG=1} prints which gate decided and why.

\code{SDPA\_BMAT\_MAX\_GB} caps the dense allocation and is enforced on \emph{every} route to
dense, so a problem that would silently need more memory than the machine has fails with a number
instead. It fails \textbf{closed}: an empty value is an error rather than ``no cap'', because
\code{SDPA\_BMAT\_MAX\_GB="\$TYPO"} is the ordinary way an empty value reaches a job script, and
treating it as no-cap would fail \emph{open} on a safety limit.

\paragraph{Available memory is reported, never acted on.} With \code{SDPA\_BMAT\_LOG=1}, and in the
message a \code{SDPA\_BMAT\_MAX\_GB} refusal prints, the dense requirement is shown next to what the
machine has free, so ``needs 21 GB'' can be read against something. Nothing branches on it, and CI
asserts that nothing does: a route chosen from free memory would be a route that changes with the
machine, which is the opposite of what this fork guarantees. The figure is Linux-only
(\code{MemAvailable} from \code{/proc/meminfo}; elsewhere the line says it is not readable rather
than guessing), and it reports the \emph{host}, so inside a memory cgroup it is an upper bound.

\subsection{Threading the sparse Cholesky}

\begin{center}\small
\begin{tabular}{@{}p{5.1cm}p{5.6cm}l@{}}
\toprule
variable & values & changes results? \\
\midrule
\code{SDPA\_SPCHOL\_MODE} & \code{auto} (default) $\cdot$ \code{serial} $\cdot$ \code{force}
  $\cdot$ \code{legacy} & no \\
\code{SDPA\_DD\_MIN\_SPCHOL\_WORK}  & integer, default \textbf{10000} & no \\
\code{SDPA\_DD\_MIN\_SPCHOL\_WIDTH} & integer, default \textbf{8} & no \\
\code{SDPA\_DD\_MIN\_SPCHOL\_TOTAL} & integer, default \textbf{0} (pass-through) & no \\
\code{SDPA\_SPCHOL\_LOG} & \code{1} on, \code{0} off & no \\
\bottomrule
\end{tabular}
\end{center}

\code{legacy} runs an independent transcription of the pre-refactor expression. It exists as an
\textbf{oracle}: \code{serial} and the threaded path share their inner kernel, so comparing those
two cannot detect a bug in the shared kernel, whereas \code{legacy} can.

\code{force} \textbf{forces or fails} --- in a build without OpenMP it is refused rather than
silently running serial, so a benchmark cannot believe it measured the threaded path when it
measured the other one.

The gates were calibrated on \code{dd}, on two architectures (i9-13900K/$24$ cores and Apple M1
Max/$10$ cores). They are runtime-overridable precisely so a machine that disagrees can be
accommodated without a rebuild.

\subsection{Threading the sparse \texorpdfstring{$\bMat$}{bMat} assembly}

\begin{center}\small
\begin{tabular}{@{}p{5.1cm}p{5.6cm}l@{}}
\toprule
variable & values & changes results? \\
\midrule
\code{SDPA\_BMAT\_ASM\_MODE} & \code{auto} (default) $\cdot$ \code{serial} $\cdot$ \code{parallel} & no \\
\code{SDPA\_BMAT\_ASM\_MIN\_PAIRS} & integer, default \textbf{8000} & no \\
\code{SDPA\_BMAT\_ASM\_SCRATCH\_MB} & integer MB, default \textbf{4096} & no \\
\code{SDPA\_BMAT\_ASM\_LOG} & \code{1} & no \\
\code{SDPA\_BMAT\_ASM\_PROFILE} & \code{1} & no \\
\code{SDPA\_BMAT\_ASM\_CENSUS} & \code{1} & no \\
\bottomrule
\end{tabular}
\end{center}

\code{SDPA\_BMAT\_ASM\_SCRATCH\_MB} is a \textbf{production safety control}, not observability: each
worker beyond the first needs two dense matrices of the largest block's order. If the budget admits
fewer workers than requested, the team shrinks; it never silently exceeds the budget.

\code{SDPA\_BMAT\_ASM\_MIN\_PAIRS = 8000} was calibrated on two architectures. Below roughly
$5{,}000$ pairs threading loses --- the fork/join costs more than the work --- and above roughly
$15{,}000$ it wins by $8$--$10\times$. The gate is a \textbf{proxy}: pair count does not fully
determine the work, since a group's cost also scales with its block dimension, so a problem near the
boundary may benefit from overriding it. A work-based replacement was attempted and the result was
\emph{negative}: no static pair-count threshold is safe on both calibration machines. That is
recorded here rather than quietly dropped.

\subsection{Profiling and census}

\code{SDPA\_BMAT\_ASM\_PROFILE=1} splits the \code{Make bMat} phase into zeroing, SDP assembly and
LP assembly. \code{SDPA\_SOLVE\_PROFILE=1} splits the triangular solve into its forward and backward
passes. Both are reported in the phase table and cost a few clock reads per iteration when off.

\code{SDPA\_BMAT\_ASM\_CENSUS=1} reports, per SDP block: pair count, group count, largest group, the
formula mix, and how many pairs consume the F2/dense-$A_j$ path. Useful for judging whether a
problem can balance across threads before trying.

\section{Exit status}

Upstream exits $0$ on \emph{every} path, including fatal errors, so a crashed run is
indistinguishable from a solved one in any harness that checks exit codes. This fork distinguishes
them.

{\small
\begin{longtable}{@{}>{\raggedright\arraybackslash}p{10.4cm}>{\raggedright\arraybackslash}p{4.6cm}@{}}
\toprule
outcome & exit \\
\midrule
\endhead
solver ran to any stopping condition (\code{pdOPT}, \code{pdFEAS}, \code{pFEAS}, \code{dFEAS},
  \code{pdINF}, \code{pINF\_dFEAS}, \code{pFEAS\_dINF}, \code{pUNBD}, \code{dUNBD},
  \code{noINFO}) & \textbf{0} \\[2pt]
iteration limit reached & \textbf{0} \\[2pt]
infeasibility / unboundedness detected & \textbf{0} \\[2pt]
malformed input, unreadable file, invalid parameter & \textbf{1}, with a diagnostic
  (line-numbered for data files) \\[2pt]
numerical failure with nothing valid to print --- no iteration completed, or the updated $X$/$Z$
  left the cone \emph{and the rollback could not be refactored} & \textbf{2},
  \code{solveStatus = FAILURE}, no solution section \\[2pt]
recoverable late failure after $k$ good iterations --- a Schur factorisation failure, or an
  $X$/$Z$ update that left the cone and \emph{was} rolled back; the last valid iterate is printed
  and labelled & \textbf{3}, \code{solveStatus = PARTIAL},
  \code{failureIteration = }$k$ \\
\bottomrule
\end{longtable}}

Infeasibility and the iteration limit are valid mathematical results, not errors, and are reported
as such.

\section{What is not established}

Stated positively so that nobody has to infer it from silence.

\begin{itemize}
\item \textbf{No time-to-solution comparison exists for the large sparse problems.} \code{dE3} and
      \code{dE4} do not converge at double-double precision under either tolerance tested
      (\code{1e-30} shipped, \code{1e-20} relaxed). Every figure quoted for them is
      \emph{per-iteration cost over a fixed four-iteration budget}. A looser parameter choice is a
      different experiment.
\item \textbf{The route census is structural, the timings are seven observations.} Seven
      representatives cover all seven classes containing the $167$ switch instances; $167$
      instances were not timed.
\item \textbf{The assembly work gate has no principled replacement yet.} The attempt to derive one
      returned a negative result on two machines.
\item \textbf{The OpenMP work thresholds carry a precision caveat.} They were calibrated at
      double-double. The sibling \code{gmp} fork inherits the mechanism but would need them
      re-measured at $256$ bits.
\item \textbf{Short cells are quoted to two significant figures on purpose.} The $24$-thread
      \code{dE4} cell spreads $8.2\%$ across nine runs in three campaigns ($6.234$--$6.761$~s), and
      the $24$-thread \code{dE3} sparse cell spreads $12\%$ across six runs. Ratios derived from
      them are ``about $7.2\times$'' and ``about $5\times$''; every earlier revision that quoted
      three figures had to move the number at the next re-measurement.
\item \textbf{Threaded upstream is nondeterministic}, so fork-versus-upstream figures are observed
      end-to-end speedups rather than strictly same-trajectory comparisons.
\end{itemize}

\section{Sources}

\begin{itemize}
\item Upstream: \url{https://github.com/nakatamaho/sdpa-dd} at \code{6eaad8d9}. The base commit is
      not an ancestor of this repository's history; fetch it explicitly to diff.
\item This fork: \url{https://github.com/Canonical111/sdpa-dd-omp}
\item Sibling fork, arbitrary precision, where the factorisation policy was derived and measured at
      full convergence: \url{https://github.com/Canonical111/sdpa-gmp-omp}
\item Benchmark tables and methodology: \code{BENCHMARKS.md}; raw per-repeat rows and per-file
      provenance under \code{bench/}, with \code{bench/validate.py} as an integrity check.
\item Build and verification: \code{INSTALL.md}.
\end{itemize}

Every modified or added source file carries an in-file, dated change notice naming its own licence:
GPLv2 \S2a for the SDPA sources, and the two-clause BSD terms for \code{mplapack/}, which carries no
GNU licence at all. The complete, always-current list of changes is a diff against upstream; an
enumeration went stale twice and is deliberately not repeated here. \code{git log} carries the
rationale per change.

\end{document}
